\section*{Chapitre \ref{ch:g12:mountain-building} --- La formation des montagnes}

\begin{solution}{exo:g12:mountain-building:1}
(a) \omterm{def:g12:mountain-building:thrust}{chevauchement}~; (b) \omterm{def:g12:mountain-building:ophiolite}{ophiolite}~; (c) \omterm{def:g12:mountain-building:fold}{synclinal}~; (d) \omterm{def:g12:mountain-building:suture}{suture}~;
(e) \omterm{def:g12:mountain-building:fold}{anticlinal}~; (f) \omterm{def:g12:mountain-building:nappe}{nappe de charriage}.
\end{solution}

\begin{solution}{exo:g12:mountain-building:2}
(1) Les flancs plongent en s'écartant de l'\omterm{def:g12:mountain-building:fold}{axe du pli}~: c'est un
\omterm{def:g12:mountain-building:fold}{anticlinal}. (2) Les couches les plus anciennes. (3) Non~: l'érosion suit
les couches dures et tendres, pas la forme du \omterm{def:g12:mountain-building:fold}{pli}~; le cœur d'un \omterm{def:g12:mountain-building:fold}{anticlinal}
peut être creusé en vallée.
\end{solution}

\begin{solution}{exo:g12:mountain-building:3}
(1) $e = 1 - \cos\ang{30} = 1 - 0.866 = 0.134$, environ \qty{13}{\percent}.
(2) $L_0 = 6.0/\cos\ang{30} \approx \qty{6.9}{km}$.
\end{solution}

\begin{solution}{exo:g12:mountain-building:4}
(1) Une \omterm{def:g11:ocean-ridges:normal-fault}{faille normale}~; la croûte a été étirée. (2) Oui~: c'est une \omterm{def:g12:mountain-building:thrust}{faille
inverse} de pendage inférieur à \ang{45}.
\end{solution}

\begin{solution}{exo:g12:mountain-building:5}
(1) Sédiments marins, basaltes en coussins, complexe filonien, \omterm{def:g11:ocean-ridges:gabbro}{gabbro},
\omterm{def:g11:ocean-ridges:peridotite}{péridotite}. (2) La lave ne forme des coussins que lorsqu'elle refroidit dans
l'eau. (3) Une \omterm{def:g12:mountain-building:ophiolite}{ophiolite} est un morceau de l'océan qui séparait deux
continents~; lors de leur collision, des écailles de son fond ont été
poussées sur leurs bords au lieu de s'enfoncer~: on les trouve donc là où
l'océan s'est refermé, le long de la \omterm{def:g12:mountain-building:suture}{suture}.
\end{solution}

\begin{solution}{exo:g12:mountain-building:6}
(1) $L_0 = 2.1 + 1.8 + 2.4 + 1.6 = \qty{7.9}{km}$~;
$e = (7.9 - 6.0)/7.9 \approx 0.24$, soit \qty{24}{\percent}. (2) $\cos\theta =
6.0/7.9 \approx 0.759$, $\theta \approx \ang{41}$.
\end{solution}

\begin{solution}{exo:g12:mountain-building:7}
(1) $45\cos\ang{10} \approx \qty{44.3}{km}$. (2) $45\sin\ang{10} \approx
\qty{7.8}{km}$. (3) De tels \omterm{def:g12:mountain-building:fold}{plis} raccourcissent de
$1 - \cos\ang{12} \approx 0.0219$~: la bande devait être large de
$44.3/0.0219 \approx \qty{2000}{km}$ avant le plissement.
\end{solution}

\begin{solution}{exo:g12:mountain-building:8}
(1) La superposition~: lue sans tenir compte de la faille, la roche du
dessus passerait pour la plus jeune. (2) $500 - 75 = 425$ millions
d'années. (3) $\qty{40000}{m} /
\qty{800}{m} = 50$ fois.
\end{solution}

\begin{solution}{exo:g12:mountain-building:9}
(1) $(30.23 - 19.38) \times 111.2 = 10.85 \times 111.2 \approx
\qty{1207}{km}$. (2) $1207/14 \approx \qty{86}{km}$ par million d'années,
soit \qty{86}{mm}, environ \qty{8.6}{cm}, par an. (3) \qty{9}{m} par siècle
font \qty{9}{cm} par an~: les deux valeurs concordent à environ
\qty{5}{\percent} près.
\end{solution}

\begin{solution}{exo:g12:mountain-building:10}
(1) $H = 35/(1 - 0.40) \approx \qty{58}{km}$. (2) $\Delta H \approx
\qty{23.3}{km}$~; elle se dresse plus haut de $23.3 \times 550/3380 \approx
\qty{3.8}{km}$, et sa \omterm{def:g12:continental-crust:root}{racine} supplémentaire est profonde de
$23.3 - 3.8 \approx \qty{19.5}{km}$.
\end{solution}

\begin{solution}{exo:g12:mountain-building:11}
(1) Le dessin multiplie $\tan\theta$ par 3~: $\tan\theta = \tan\ang{40}/3
\approx 0.280$, $\theta \approx \ang{15.6}$. (2) Pendage vrai~: $1 - \cos\ang{15.6}
\approx 0.037$, soit \qty{3.7}{\percent}~; pendage lu sur le dessin~: $1 - \cos\ang{40}
\approx 0.234$, soit \qty{23}{\percent}. (3) Une échelle verticale étirée
redresse les pendages et allonge les couches~: elle fait paraître le
\omterm{def:g12:mountain-building:shortening}{raccourcissement} plusieurs fois trop grand.
\end{solution}

\begin{solution}{exo:g12:mountain-building:12}
(1) $L_0 = 10/(1 - 0.317) \approx \qty{14.6}{km}$. (2) $1 - \cos\ang{57.5}
\approx 1 - 0.537 = 0.463$, soit \qty{46}{\percent}. (3) Dans le \omterm{def:g12:mountain-building:fold}{pli}
arrondi, seule une courte partie de la couche plonge de \ang{57.5}~; près
des charnières, elle est presque plate, et les parties plates ne sont pas
raccourcies. Dans le \omterm{def:g12:mountain-building:fold}{pli} à flancs droits, toute la couche plonge de
\ang{57.5}.
\end{solution}

\begin{solution}{exo:g12:mountain-building:13}
(1) $\sqrt{16.6^2 + 41.8^2} \approx \qty{45.0}{mm/yr}$. (2) $41.8 - 23.5 =
\qty{18.3}{mm}$ par an, soit environ \qty{18}{km} de \omterm{def:g12:mountain-building:shortening}{raccourcissement}
nord--sud par million d'années au nord de Lhassa. (3) Lhassa elle-même
s'est déplacée vers le nord, et les séismes de la ceinture s'étendent loin
au nord de l'Himalaya, à travers l'Asie~: la croûte continue de s'y
rompre.
\end{solution}

\begin{solution}{exo:g12:mountain-building:14}
(1) Son pendage est presque nul~: le \omterm{def:g12:mountain-building:shortening}{raccourcissement} est égal au
déplacement, \qty{45}{km}. (2) $45/(1 - \cos\ang{12}) = 45/0.0219 \approx \qty{2060}{km}$.
(3) Les grandes convergences sont absorbées surtout par les \omterm{def:g12:mountain-building:thrust}{chevauchements}
et les \omterm{def:g12:mountain-building:nappe}{nappes de charriage}, et par les \omterm{def:g12:mountain-building:fold}{plis} serrés~; les \omterm{def:g12:mountain-building:fold}{plis} doux
absorbent très peu.
\end{solution}

\begin{solution}{exo:g12:mountain-building:15}
(1) $\qty{8848.86}{m} = \qty{8848860}{mm}$ en $\num{5e7}$ ans~: environ
\qty{0.18}{mm} par an. (2) $40/0.18 \approx 220$~: la montée est plus de
deux cents fois plus lente que le mouvement de l'Inde. (3) Le calcaire
s'est déposé sous le niveau de la mer, et l'érosion a enlevé les roches qui
le recouvraient~: il a été soulevé de plus que son altitude actuelle.
\end{solution}

\begin{solution}{pb:g12:mountain-building:1}
\textbf{1.} $2\pi \times 6371/360 \approx \qty{111.2}{km}$.

\textbf{2.} Entre il y a 50 millions d'années (\ang{0.68}~S) et il y a 45
(\ang{3.30}~N)~: il y a $50 - 5 \times 0.68/3.98 \approx 49$ millions
d'années.

\textbf{3.} $(19.38 - 0.68) \times 111.2 = 18.70 \times 111.2 \approx
\qty{2079}{km}$.

\textbf{4.} $2079/16 \approx \qty{130}{km}$ par million d'années~:
\qty{130}{mm} par an, environ \qty{13}{cm}.

\textbf{5.} De 66 à 60~: $7.28 \times 111.2/6 \approx 135$~; de 60 à 55~: $4.86
\times 111.2/5 \approx 108$~; de 55 à 50~: $6.56 \times 111.2/5 \approx 146$~; de 50
à 45~: $3.98 \times 111.2/5 \approx 89$~; de 45 à 40~: $2.29 \times 111.2/5
\approx 51$ (\unit{mm} par an). La plus rapide entre 55 et 50 millions
d'années.

\textbf{6.} $(20.00 - 5.59) \times 111.2/40 = 14.41 \times 111.2/40 \approx
\qty{40}{mm}$ par an.

\textbf{7.} $40/130 \approx 0.31$~: moins d'un tiers.

\textbf{8.} De l'étape 55--50 (146) à l'étape 50--45 (89)~: vers 50
millions d'années.

\textbf{9.} La \omterm{def:g12:continental-crust:continental-crust}{croûte continentale} est trop légère pour s'enfoncer~: quand
elle a atteint la fosse, la \omterm{def:g11:subduction-zones:subduction}{subduction} n'a plus pu entraîner librement
l'Inde, et l'autre continent l'a freinée.

\textbf{10.} $(20.00 + 0.68) \times 111.2 = 20.68 \times 111.2 \approx
\qty{2300}{km}$.

\textbf{11.} $(29.65 - 21.15) \times 111.2 = 8.50 \times 111.2 \approx
\qty{945}{km}$.

\textbf{12.} $2300 - 945 \approx \qty{1354}{km}$.

\textbf{13.} Aujourd'hui, $(29.65 - 20.00) \times 111.2 \approx \qty{1073}{km}$~;
il y a 50 millions d'années, $(21.15 + 0.68) \times 111.2 = 21.83 \times 111.2
\approx \qty{2427}{km}$ (et $2427 - 1073 = 1354$).

\textbf{14.} $e = 1354/2427 \approx 0.56$, soit \qty{56}{\percent}.

\textbf{15.} Les \omterm{def:g12:mountain-building:fold}{plis}, et les \omterm{def:g12:mountain-building:thrust}{failles inverses}~: des \omterm{def:g12:mountain-building:thrust}{chevauchements} qui
transportent des \omterm{def:g12:mountain-building:nappe}{nappes de charriage}.

\textbf{16.} $H = 35 \times 2427/1073 \approx \qty{79}{km}$.

\textbf{17.} $\Delta H = 79 - 35 = \qty{44}{km}$~; elle se dresserait plus
haut de $44 \times 550/3380 \approx \qty{7.2}{km}$.

\textbf{18.} $\Delta H = 4.6 \times 3380/550 \approx \qty{28}{km}$, soit une
croûte épaisse de $35 + 28 = \qty{63}{km}$, proche des \qty{70}{km}
mesurés.

\textbf{19.} La croûte ne garde pas sa section~: des parties de l'Asie sont
chassées vers l'est, et l'érosion enlève de la roche au sommet des
montagnes.

\textbf{20.} $2079 + 2300 \approx \qty{4380}{km}$~: depuis la fin du
Crétacé, l'Inde a parcouru environ \textbf{\qty{4400}{km} vers le nord}.
\end{solution}
